% One response, drawn as its tree of positions: a null raised at the speaker
% of session index 3 climbs until it reaches a position whose type is allowed
% to hold it.
%
% Same stroke weight, arrowhead and font as figures/tikz/ch01-one-field.tex,
% the file that fixes this book's idiom: 0.4pt strokes, the Straight Barb
% arrowhead, sans-serif scriptsize labels, square corners, no fills, black on
% white. That figure and ch03-one-batch.tex are timelines because their
% subject is order. This one is a shape, so it is a tree; the idiom is the
% line weight and the labelling, not the axis.
%
% Solid box: a position that is present in the response. That includes nodes,
% whose key survives even though its value is null. Dashed box: a position
% that does not exist in the response at all, because the list that would
% have held it is gone. Each box on the path carries its own GraphQL type,
% because the nullability of that type is the only thing deciding whether the
% climb stops there.
%
% The two solid arrows are the climb, drawn along the tree edges they run up.
% The query's title and name fields are left out: neither changes where the
% climb stops.
%
% Bare tikzpicture (house style): the figure environment, caption and label
% live at the call site in the section file.
\begin{tikzpicture}[
  x=1mm, y=1mm,
  every node/.append style={font=\sffamily\scriptsize},
  posbox/.style={draw, line width=0.4pt, minimum width=20mm,
    minimum height=9mm, align=center, inner sep=1pt},
  lostbox/.style={posbox, dashed},
  link/.style={draw, line width=0.4pt},
  lostlink/.style={draw, line width=0.4pt, dashed},
  climb/.style={draw, line width=0.4pt,
    -{Straight Barb[length=1.4mm, width=1.6mm]}},
  span/.style={draw, line width=0.4pt},
  spanlabel/.style={anchor=north, align=center, inner sep=2pt},
  note/.style={anchor=west, align=left, inner sep=2pt},
  haltnote/.style={anchor=east, align=right, text width=30mm, inner sep=2pt},
]

% --------------------------------------------------------------- positions
\node[posbox] (root)  at (83,78) {sessions};

\node[posbox] (nodes) at (46,54) {nodes\\{}[Session!]};
\node[posbox] (page)  at (120,54) {pageInfo};

\node[lostbox] (i0) at (10,30) {nodes[0]};
\node[lostbox] (i1) at (32,30) {nodes[1]};
\node[lostbox] (i2) at (54,30) {nodes[2]};
\node[lostbox] (i3) at (82,30) {nodes[3]\\Session!};

\node[lostbox] (speaker) at (82,6)  {speaker\\Speaker!};
\node[posbox]  (hasnext) at (120,30) {hasNextPage};

% ---------------------------------------------------------- tree structure
\draw[link] (root.south) -- (83,66) -- (46,66) -- (nodes.north);
\draw[link] (root.south) -- (83,66) -- (120,66) -- (page.north);
\draw[link] (page.south) -- (hasnext.north);

\draw[lostlink] (nodes.south) -- (46,42) -- (10,42) -- (i0.north);
\draw[lostlink] (nodes.south) -- (46,42) -- (32,42) -- (i1.north);
\draw[lostlink] (nodes.south) -- (46,42) -- (54,42) -- (i2.north);

% ------------------------------------------------------------- the climb
\draw[climb] (speaker.north) -- (i3.south);
\draw[climb] (i3.north) -- (82,42) -- (46,42) -- (nodes.south);

% ------------------------------------------------------------------ notes
% Each note sits on a row of its own. The types are carried inside the boxes
% rather than in side labels, because a label beside nodes[3] runs into the
% pageInfo branch at this width.
\node[note] at (94,6)
  {the resolver returned null here,\\and Speaker! cannot hold it};
\node[haltnote] at (32,54)
  {the first position that may be null, so the climb stops and nodes is
   null};
\node[spanlabel] at (120,24) {untouched};

% ------------------------------------------------------- collateral loss
\draw[span] (0,23) -- (0,19) -- (64,19) -- (64,23);
\node[spanlabel] at (32,17)
  {these three resolved correctly and are gone anyway};

% ----------------------------------------------------------------- legend
\node[posbox, minimum width=7mm, minimum height=4mm] at (4,-2) {};
\node[note] at (9,-2) {present in the response};

\node[lostbox, minimum width=7mm, minimum height=4mm] at (62,-2) {};
\node[note] at (67,-2) {not in the response at all};

\end{tikzpicture}
